\documentclass[11pt,a4paper]{article}

\usepackage[utf8]{inputenc}
\usepackage[T1]{fontenc}
\usepackage[provide=*,english]{babel}
\usepackage{lmodern}
\usepackage{microtype}
\usepackage{amsmath,amssymb,mathtools}
\usepackage{geometry}
\usepackage{xcolor}
\usepackage{enumitem}
\usepackage{booktabs}
\usepackage{fancyhdr}
\usepackage[hidelinks]{hyperref}

\geometry{left=2.3cm,right=2.3cm,top=2.2cm,bottom=2.25cm,headheight=14pt}
\setlength{\parindent}{0pt}
\setlength{\parskip}{0.38em}
\setlist[itemize]{topsep=0.25em,itemsep=0.18em,leftmargin=1.5em}
\allowdisplaybreaks

\definecolor{blue}{RGB}{34,73,110}
\definecolor{gray}{RGB}{90,96,102}

\pagestyle{fancy}
\fancyhf{}
\fancyhead[L]{\small\color{gray}GC2D: Gauss--Legendre method}
\fancyhead[R]{\small\color{gray}Jacobian, symplecticity, and symmetry}
\fancyfoot[C]{\thepage}
\renewcommand{\headrulewidth}{0.3pt}

\newcommand{\R}{\mathbb{R}}
\newcommand{\Id}{I}
\newcommand{\Om}{\Omega}
\newcommand{\Jn}{\mathcal{J}_n}
\newcommand{\Phih}{\Phi_{h,t_n}}

\begin{document}

\begin{center}
  {\LARGE\bfseries\color{blue}Two-stage Gauss--Legendre Method}\par
  \vspace{0.35em}
  {\large Exact one-step Jacobian, symplecticity, and symmetry}\par
  \vspace{0.45em}
  {\small Fourth-order collocation method for canonical Hamiltonian systems}
\end{center}

\vspace{0.5em}

\section{Hamiltonian convention}

Let $z\in\R^{2d}$ and consider the possibly time-dependent canonical
Hamiltonian system
\begin{equation}
 \dot z=f(t,z)=\Om^{-1}\nabla_z H(t,z),
 \qquad
 \Om=
 \begin{pmatrix}
  0&\Id_d\\
  -\Id_d&0
 \end{pmatrix}.
 \label{eq:hamiltonian-system}
\end{equation}
For the component-major guiding-centre state
\[
 z=(x_1,\ldots,x_d,y_1,\ldots,y_d)^T,
\]
this convention gives
\[
 f(t,z)=
 \begin{pmatrix}
  -\nabla_y H(t,z)\\
  \nabla_x H(t,z)
 \end{pmatrix}.
\]
At any fixed time and state, let $F=D_zf(t,z)$.  Since
$\nabla_z^2H$ is symmetric,
\begin{equation}
 F^T\Om+\Om F=0.
 \label{eq:hamiltonian-jacobian}
\end{equation}
This identity is the infinitesimal Hamiltonian condition used below.  The
proof remains valid for a time-dependent Hamiltonian because the stage times
do not depend on the initial state $z_n$.

\section{Gauss--Legendre stages}

Set
\[
 r=\frac{\sqrt{3}}{6},
 \qquad
 c_1=\frac12-r,
 \qquad
 c_2=\frac12+r.
\]
The Butcher coefficients of the two-stage Gauss--Legendre method are
\begin{equation}
 A=
 \begin{pmatrix}
  \frac14&\frac14-r\\[1mm]
  \frac14+r&\frac14
 \end{pmatrix},
 \qquad
 b=
 \begin{pmatrix}
  \frac12\\[1mm]
  \frac12
 \end{pmatrix}.
 \label{eq:tableau}
\end{equation}
For one step of length $h$, the stage states satisfy
\begin{equation}
 Z_i=z_n+h\sum_{j=1}^2a_{ij}f(t_n+c_jh,Z_j),
 \qquad i=1,2,
 \label{eq:stage-equations}
\end{equation}
and the accepted state is
\begin{equation}
 z_{n+1}=\Phih(z_n)
 =z_n+h\sum_{i=1}^2b_i f(t_n+c_ih,Z_i).
 \label{eq:step-map}
\end{equation}
The method is implicit: equations \eqref{eq:stage-equations} determine
$(Z_1,Z_2)$ simultaneously.

\clearpage

\section{Direct proof of fourth-order accuracy}

We now prove directly why the Runge--Kutta order conditions through degree four
give a fourth-order method.  For clarity, first consider the autonomous
problem
\begin{equation}
 \dot y=f(y),
 \qquad y(t_n)=y.
 \label{eq:autonomous-order-problem}
\end{equation}
The non-autonomous case follows by augmenting the state with time,
$\widehat y=(t,y)$ and $\dot{\widehat y}=(1,f(t,y))$.  Because
$c=A\mathbf{1}$, applying the Runge--Kutta method to this augmented autonomous
system produces exactly the stage times used in
\eqref{eq:stage-equations}.

Set
\begin{equation}
 F=f(y),
 \qquad
 F'=Df(y),
 \qquad
 F''=D^2f(y),
 \qquad
 F'''=D^3f(y),
 \label{eq:elementary-differential-notation}
\end{equation}
where the derivatives are multilinear maps.  Repeated application of the
chain rule gives
\begin{align}
 y'&=F,
 \notag\\
 y''&=F'F,
 \notag\\
 y'''&=F''(F,F)+F'^2F,
 \notag\\
 y^{(4)}&=F'''(F,F,F)
 +3F''(F,F'F)
 +F'F''(F,F)
 +F'^3F.
 \label{eq:exact-time-derivatives}
\end{align}
Consequently, the exact flow has the expansion
\begin{align}
 y(t_n+h)={}&y+hF+\frac{h^2}{2}F'F
 \notag\\
 &+\frac{h^3}{6}
 \left[F''(F,F)+F'^2F\right]
 \notag\\
 &+\frac{h^4}{24}
 \left[
 F'''(F,F,F)
 +3F''(F,F'F)
 +F'F''(F,F)
 +F'^3F
 \right]
 +O(h^5).
 \label{eq:exact-order-four-expansion}
\end{align}

For a generic Runge--Kutta step, write
\begin{equation}
 K_i=f(Y_i),
 \qquad
 Y_i=y+h\sum_j a_{ij}K_j,
 \qquad
 y_{n+1}=y+h\sum_i b_iK_i,
 \label{eq:generic-rk-order-step}
\end{equation}
and define $C=\operatorname{diag}(c)$.  Powers such as $c^2$ and $c^3$
are understood componentwise.  If
$\Delta_i=h\sum_j a_{ij}K_j$, Taylor expansion about $y$ gives
\begin{equation}
 K_i
 =F+F'\Delta_i
 +\frac12F''(\Delta_i,\Delta_i)
 +\frac16F'''(\Delta_i,\Delta_i,\Delta_i)
 +O(h^4).
 \label{eq:stage-field-taylor}
\end{equation}
Substituting the stage expansions recursively into
\eqref{eq:stage-field-taylor} and collecting equal powers of $h$ yields
\begin{align}
 K_i={}&F+h c_iF'F
 \notag\\
 &+h^2\left[
 (Ac)_iF'^2F
 +\frac12c_i^2F''(F,F)
 \right]
 \notag\\
 &+h^3\left[
 (A^2c)_iF'^3F
 +\frac12(Ac^2)_iF'F''(F,F)
 \right.
 \notag\\
 &\hspace{27mm}\left.
 +c_i(Ac)_iF''(F,F'F)
 +\frac16c_i^3F'''(F,F,F)
 \right]
 +O(h^4).
 \label{eq:rk-stage-order-four-expansion}
\end{align}
Inserting \eqref{eq:rk-stage-order-four-expansion} into the accepted update
in \eqref{eq:generic-rk-order-step} gives
\begin{align}
 y_{n+1}={}&y+h(b^T\mathbf{1})F
 +h^2(b^Tc)F'F
 \notag\\
 &+h^3\left[
 (b^TAc)F'^2F
 +\frac12(b^Tc^2)F''(F,F)
 \right]
 \notag\\
 &+h^4\left[
 (b^TA^2c)F'^3F
 +\frac12(b^TAc^2)F'F''(F,F)
 \right.
 \notag\\
 &\hspace{25mm}\left.
 +(b^TCAc)F''(F,F'F)
 +\frac16(b^Tc^3)F'''(F,F,F)
 \right]
 +O(h^5).
 \label{eq:rk-order-four-expansion}
\end{align}

The elementary differentials in
\eqref{eq:exact-order-four-expansion} and
\eqref{eq:rk-order-four-expansion} are independent for a general smooth
vector field.  Their coefficients agree through $h^4$ precisely when
\begin{align}
 b^T\mathbf{1}&=1,
 &b^Tc&=\frac12,
 \notag\\
 b^Tc^2&=\frac13,
 &b^TAc&=\frac16,
 \notag\\
 b^Tc^3&=\frac14,
 &b^TCAc&=\frac18,
 \notag\\
 b^TAc^2&=\frac1{12},
 &b^TA^2c&=\frac1{24}.
 \label{eq:order-four-conditions}
\end{align}
For example, the coefficient of $F''(F,F)$ in the numerical step is
$\tfrac12b^Tc^2$ and its exact coefficient is $\tfrac16$, which gives
$b^Tc^2=\tfrac13$.  The other seven equalities follow from the same
term-by-term comparison.  Thus these conditions are not merely formal
identities: they are exactly the equations that make the numerical and exact
Taylor expansions coincide through fourth degree.

It remains to verify them for the coefficients in \eqref{eq:tableau}.  Direct
calculation gives
\begin{equation}
 b^T\mathbf{1}=1,
 \qquad
 b^Tc=\frac12,
 \qquad
 b^Tc^2=\frac13,
 \qquad
 b^Tc^3=\frac14,
 \label{eq:gauss-moment-identities}
\end{equation}
and
\begin{equation}
 Ac=\frac12c^2,
 \qquad
 b^TA=b^T(I-C).
 \label{eq:gauss-stage-order-identities}
\end{equation}
Therefore
\begin{align}
 b^TAc&=\frac12b^Tc^2=\frac16,
 &b^TCAc&=\frac12b^Tc^3=\frac18,
 \notag\\
 b^TAc^2&=b^T(c^2-c^3)=\frac1{12},
 &b^TA^2c&=\frac12b^TAc^2=\frac1{24}.
 \label{eq:gauss-remaining-order-identities}
\end{align}
Hence all conditions in \eqref{eq:order-four-conditions} hold, and comparison
of \eqref{eq:exact-order-four-expansion} with
\eqref{eq:rk-order-four-expansion} proves the one-step defect estimate
\begin{equation}
 \left\lVert y(t_n+h)-y_{n+1}\right\rVert=O(h^5).
 \label{eq:local-order-five-defect}
\end{equation}

Finally, let $e_n$ be the global error.  For a smooth Lipschitz vector field,
the Runge--Kutta one-step map is Lipschitz with factor $1+Lh$ for sufficiently
small $h$.  The local estimate \eqref{eq:local-order-five-defect} then gives
\begin{equation}
 \lVert e_{n+1}\rVert
 \leq(1+Lh)\lVert e_n\rVert+Ch^5.
 \label{eq:global-error-recurrence}
\end{equation}
For $nh\leq T$, the discrete Gronwall estimate yields
\begin{equation}
 \lVert e_n\rVert
 \leq Ch^5\sum_{j=0}^{n-1}(1+Lh)^j
 \leq C_T h^4.
 \label{eq:global-order-four-bound}
\end{equation}
Thus the two-stage Gauss--Legendre method has global order four.

\section{Exact Jacobian of the stage root}

Introduce the stage field Jacobians and stage sensitivities
\begin{equation}
 F_i:=D_zf(t_n+c_ih,Z_i),
 \qquad
 S_i:=\frac{\partial Z_i}{\partial z_n}.
 \label{eq:stage-derivatives}
\end{equation}
Differentiating \eqref{eq:stage-equations} with respect to $z_n$ gives
\begin{equation}
 S_i=\Id_{2d}+h\sum_{j=1}^2a_{ij}F_jS_j.
 \label{eq:stage-sensitivities}
\end{equation}
Equivalently,
\begin{equation}
 \underbrace{
 \begin{pmatrix}
  \Id_{2d}-\frac h4F_1
   &-h\left(\frac14-r\right)F_2\\[1mm]
  -h\left(\frac14+r\right)F_1
   &\Id_{2d}-\frac h4F_2
 \end{pmatrix}}_{M=D_{(Z_1,Z_2)}R}
 \begin{pmatrix}S_1\\S_2\end{pmatrix}
 =
 \begin{pmatrix}\Id_{2d}\\\Id_{2d}\end{pmatrix}.
 \label{eq:stage-tangent-system}
\end{equation}
Here $M$ is the Jacobian of the coupled stage residual.  It is also the matrix
used by an exact Newton linearization, but it is \emph{not} the Jacobian of the
physical one-step map.

Assuming that the converged stage root is regular, the implicit function
theorem gives
\begin{equation}
 \begin{pmatrix}S_1\\S_2\end{pmatrix}
 =M^{-1}
 \begin{pmatrix}\Id_{2d}\\\Id_{2d}\end{pmatrix}.
 \label{eq:stage-tangent-solution}
\end{equation}
In an implementation, equation \eqref{eq:stage-tangent-system} should be
solved directly; the inverse in \eqref{eq:stage-tangent-solution} is only
notation.

\section{Exact Jacobian of the physical step}

Define the Jacobian of the map \eqref{eq:step-map} by
\[
 \Jn:=D_{z_n}\Phih(z_n).
\]
Differentiating the accepted update and using
\eqref{eq:stage-derivatives} yields
\begin{equation}
 \boxed{
 \Jn=\Id_{2d}+h\sum_{i=1}^2b_iF_iS_i
 =\Id_{2d}+\frac h2\left(F_1S_1+F_2S_2\right).
 }
 \label{eq:step-jacobian}
\end{equation}
Combining \eqref{eq:stage-tangent-solution} and
\eqref{eq:step-jacobian}, the same Jacobian can be written as
\begin{equation}
 \boxed{
 \Jn
 =\Id_{2d}
 +h
 \begin{pmatrix}\frac12F_1&\frac12F_2\end{pmatrix}
 M^{-1}
 \begin{pmatrix}\Id_{2d}\\\Id_{2d}\end{pmatrix}.
 }
 \label{eq:explicit-step-jacobian}
\end{equation}
This is the exact tangent of the ideal converged root.  It differentiates the
mathematical Gauss--Legendre map, not a finite Newton stopping rule.

\clearpage

\section{Direct proof of symplecticity with the step Jacobian}

The one-step map is symplectic if and only if
\begin{equation}
 \Jn^T\Om\Jn=\Om.
 \label{eq:symplectic-condition}
\end{equation}
Set
\begin{equation}
 K_i:=F_iS_i.
 \label{eq:Ki}
\end{equation}
Then equations \eqref{eq:step-jacobian} and
\eqref{eq:stage-sensitivities} become
\begin{equation}
 \Jn=\Id_{2d}+h\sum_i b_iK_i,
 \qquad
 \Id_{2d}=S_i-h\sum_j a_{ij}K_j.
 \label{eq:compact-identities}
\end{equation}
Expanding the symplecticity defect gives
\begin{align}
 \Jn^T\Om\Jn-\Om
 ={}&h\sum_i b_i\left(K_i^T\Om+\Om K_i\right)
 \notag\\
 &+h^2\sum_{i,j}b_ib_jK_i^T\Om K_j.
 \label{eq:defect-expansion}
\end{align}
By the Hamiltonian identity \eqref{eq:hamiltonian-jacobian},
\begin{align}
 K_i^T\Om S_i+S_i^T\Om K_i
 &=S_i^T\left(F_i^T\Om+\Om F_i\right)S_i
 \notag\\
 &=0.
 \label{eq:hamiltonian-cancellation}
\end{align}
Using the second identity in \eqref{eq:compact-identities} on both occurrences
of $\Id_{2d}$ therefore gives
\begin{equation}
 K_i^T\Om+\Om K_i
 =-h\sum_j a_{ij}
 \left(K_i^T\Om K_j+K_j^T\Om K_i\right).
 \label{eq:linear-term-reduction}
\end{equation}
Substitute \eqref{eq:linear-term-reduction} into
\eqref{eq:defect-expansion}.  Renaming $i$ and $j$ in the second term arising
from \eqref{eq:linear-term-reduction} produces
\begin{equation}
 \boxed{
 \Jn^T\Om\Jn-\Om
 =h^2\sum_{i,j}
 \left(b_ib_j-b_i a_{ij}-b_j a_{ji}\right)
 K_i^T\Om K_j.
 }
 \label{eq:rk-symplecticity-defect}
\end{equation}
Consequently, the defect vanishes whenever the Runge--Kutta coefficients obey
\begin{equation}
 \boxed{
 b_i a_{ij}+b_j a_{ji}=b_ib_j,
 \qquad i,j=1,2.
 }
 \label{eq:rk-symplecticity-condition}
\end{equation}

\clearpage

\section{Verification for the Gauss--Legendre coefficients}

For the diagonal terms in \eqref{eq:rk-symplecticity-condition},
\begin{equation}
 2b_1a_{11}-b_1^2
 =2\left(\frac12\right)\left(\frac14\right)-\frac14=0,
 \qquad
 2b_2a_{22}-b_2^2=0.
 \label{eq:diagonal-check}
\end{equation}
For the off-diagonal terms,
\begin{align}
 b_1a_{12}+b_2a_{21}-b_1b_2
 &=\frac12\left(\frac14-r\right)
   +\frac12\left(\frac14+r\right)-\frac14
 \notag\\
 &=0.
 \label{eq:off-diagonal-check}
\end{align}
The case $(i,j)=(2,1)$ is the same equality with the two summands exchanged.
Thus every coefficient multiplying $K_i^T\Om K_j$ in
\eqref{eq:rk-symplecticity-defect} is zero, and hence
\begin{equation}
 \boxed{\Jn^T\Om\Jn=\Om.}
 \label{eq:final-symplecticity}
\end{equation}
Therefore the two-stage Gauss--Legendre method is symplectic when its implicit
stage equations are solved exactly.

\clearpage

\section{Symmetry and time reversibility}

Symmetry is a property distinct from symplecticity.  The method is symmetric
when its backward step is exactly the inverse of its forward step.  For the
possibly time-dependent problem \eqref{eq:hamiltonian-system}, this means
\begin{equation}
 \boxed{
 \Phi_{-h,t_n+h}\circ\Phi_{h,t_n}=\operatorname{Id}.
 }
 \label{eq:symmetry-definition}
\end{equation}
For an autonomous problem, equation \eqref{eq:symmetry-definition} reduces to
$\Phi_{-h}=\Phi_h^{-1}$.

For the forward step, write
\begin{equation}
 f_1=f(t_n+c_1h,Z_1),
 \qquad
 f_2=f(t_n+c_2h,Z_2).
 \label{eq:forward-stage-fields}
\end{equation}
Equations \eqref{eq:stage-equations} and \eqref{eq:step-map} are then
\begin{align}
 Z_1&=z_n+h(a_{11}f_1+a_{12}f_2),
 \notag\\
 Z_2&=z_n+h(a_{21}f_1+a_{22}f_2),
 \label{eq:forward-stages-expanded}
\\
 z_{n+1}&=z_n+h(b_1f_1+b_2f_2).
 \label{eq:forward-output-expanded}
\end{align}

Now start from $z_{n+1}$ at time $t_n+h$ and apply the same method with step
$-h$.  Its stages $\widehat Z_1,\widehat Z_2$ satisfy
\begin{equation}
 \widehat Z_i=z_{n+1}
 -h\sum_{j=1}^2a_{ij}
 f(t_n+h-c_jh,\widehat Z_j).
 \label{eq:backward-stage-equations}
\end{equation}
The Gauss nodes are reflected about the midpoint:
\begin{equation}
 1-c_1=c_2,
 \qquad
 1-c_2=c_1.
 \label{eq:reflected-nodes}
\end{equation}
Thus the backward stage times are the forward stage times in reverse order.
We shall verify that the backward stages are
\begin{equation}
 \boxed{
 \widehat Z_1=Z_2,
 \qquad
 \widehat Z_2=Z_1.
 }
 \label{eq:reversed-stages}
\end{equation}

Under the identification \eqref{eq:reversed-stages}, the first backward stage
equation becomes
\begin{align}
 \widehat Z_1
 &=z_{n+1}-h(a_{11}f_2+a_{12}f_1)
 \notag\\
 &=z_n+h\left[(b_1-a_{12})f_1+(b_2-a_{11})f_2\right].
 \label{eq:first-backward-stage}
\end{align}
For the coefficients \eqref{eq:tableau},
\begin{equation}
 b_1-a_{12}=a_{21},
 \qquad
 b_2-a_{11}=a_{22}.
 \label{eq:first-reversal-identities}
\end{equation}
Consequently,
\[
 \widehat Z_1=z_n+h(a_{21}f_1+a_{22}f_2)=Z_2.
\]
Likewise, the second backward stage satisfies
\begin{align}
 \widehat Z_2
 &=z_{n+1}-h(a_{21}f_2+a_{22}f_1)
 \notag\\
 &=z_n+h\left[(b_1-a_{22})f_1+(b_2-a_{21})f_2\right].
 \label{eq:second-backward-stage}
\end{align}
Since
\begin{equation}
 b_1-a_{22}=a_{11},
 \qquad
 b_2-a_{21}=a_{12},
 \label{eq:second-reversal-identities}
\end{equation}
we obtain
\[
 \widehat Z_2=z_n+h(a_{11}f_1+a_{12}f_2)=Z_1.
\]
The proposed reversed stages therefore solve the coupled backward stage
equations exactly.

\clearpage

The output of the backward step is now
\begin{align}
 \widehat z_n
 &=z_{n+1}-h(b_1f_2+b_2f_1)
 \notag\\
 &=z_n+h(b_1f_1+b_2f_2)-h(b_1f_2+b_2f_1)
 \notag\\
 &=z_n,
 \label{eq:backward-recovers-input}
\end{align}
where the last equality uses $b_1=b_2=1/2$.  This proves
\eqref{eq:symmetry-definition}.

More generally, an $s$-stage Runge--Kutta method is symmetric when its
coefficients are invariant under stage reversal:
\begin{equation}
 \boxed{
 c_i=1-c_{s+1-i},
 \qquad
 b_i=b_{s+1-i},
 \qquad
 a_{ij}=b_{s+1-j}-a_{s+1-i,s+1-j}.
 }
 \label{eq:rk-symmetry-condition}
\end{equation}
The coefficients in \eqref{eq:tableau} satisfy all three identities for
$s=2$.

Differentiating the map identity \eqref{eq:symmetry-definition} also gives the
Jacobian form of reversibility:
\begin{equation}
 \boxed{
 D\Phi_{-h,t_n+h}(z_{n+1})\,
 D\Phi_{h,t_n}(z_n)=\Id_{2d}.
 }
 \label{eq:jacobian-reversibility}
\end{equation}
Thus the backward-step Jacobian is exactly the inverse of the forward-step
Jacobian.  As with symplecticity, a finite Newton tolerance and floating-point
round-off can produce a small measured reversibility defect even though the
ideal implicit method is exactly symmetric.

\clearpage

\section{Implemented configurations}

The public \texttt{GaussLegendre4} model uses the same collocation equations
for every configuration.  The selector
\texttt{newton\_jacobian\_method} accepts \texttt{analytic},
\texttt{finite\_difference}, or \texttt{auto}.  Analytic mode uses exact
guiding-centre particle blocks, finite-difference mode differentiates the
coupled stage residual, and auto resolves the available dynamics capability.
Absolute and relative tolerances and the maximum iteration count control the
finite root solve.

With \texttt{track\_energy=True}, the method also advances the time-conjugate
momentum using the two converged Gauss stages and monitors the generalized
Hamiltonian $K=H+k$.  This triangular extension does not alter the physical
collocation stages.  These are configurations of one Gauss--Legendre model,
not separate numerical methods.

\section{Consequences and numerical interpretation}

For one guiding-centre particle, $\Jn\in\R^{2\times2}$ and
\[
 \Jn^T\Om\Jn=\det(\Jn)\Om.
\]
Equation \eqref{eq:final-symplecticity} therefore implies
\[
 \det(\Jn)=1,
\]
so the one-step map preserves oriented area.  In dimensions greater than two,
unit determinant is only a consequence of symplecticity, not an equivalent
test.

If $\Jn$ denotes the local Jacobian at step $n$, the accumulated flow tangent
after $N$ steps is
\[
 \mathcal{P}_N=\mathcal{J}_{N-1}\cdots\mathcal{J}_1\mathcal{J}_0.
\]
Since a product of symplectic matrices is symplectic,
\[
 \mathcal{P}_N^T\Om\mathcal{P}_N=\Om.
\]

The algebraic result concerns the exact root and the exact tangent
\eqref{eq:step-jacobian}.  In floating-point computations, a finite Newton
tolerance, round-off, or a finite-difference approximation of the map
Jacobian produces a small nonzero measured defect
\[
 \left\lVert\Jn^T\Om\Jn-\Om\right\rVert.
\]
That numerical floor does not contradict the exact symplecticity proved above.

\section*{References}

\begin{itemize}
 \item J.~M. Sanz-Serna, ``Runge--Kutta schemes for Hamiltonian systems,''
       \emph{BIT Numerical Mathematics}, 28 (1988), 877--883.
 \item E. Hairer, C. Lubich, and G. Wanner,
       \emph{Geometric Numerical Integration}, second edition,
       Springer, 2006.
\end{itemize}

\end{document}
